% \section{Related Work}
% %merged version:


% \noindent\textbf{Reinforcement Learning for Reasoning in LLMs.}~~
% Reinforcement learning (RL) on LLMs \citep{christiano2023deepreinforcementlearninghuman,ouyang2022traininglanguagemodelsfollow,chen2021decisiontransformer} has significantly improved LLMs' reasoning capabilities. Notable approaches include the use of Proximal Policy Optimization Algorithms (PPO) \citep{schulman2017proximalpolicyoptimizationalgorithms} which maintains training stability while enhancing performance by clipping policy updates, Group Relative Policy Optimization (GRPO) \citep{deepseekai2025deepseekr1incentivizingreasoningcapability} for enhancing the ability of systematic problem-solving, actor-critic methods like SAC \citep{haarnoja2018softactorcritic} and ArCHer \citep{archer} that leverages an critic to promote robust exploration and stability, and meta tokens \citep{goyal2024thinkspeaktraininglanguage,herel2024thinkingtokenslanguagemodeling,pfau2024letsthinkdotdot} for structured thinking. Other significant developments include Process Reward Model (PRM) \citep{zhang2025lessonsdevelopingprocessreward, lightman2023verify} and Monte Carlo Tree Search (MCTS) based approaches \citep{hao2023reasoning} for systematic problem-solving. On the other hand, recent advances in LLM reasoning have explored techniques to enable models to generate intermediate chain-of-thought rationales. In particular, STaR \citep{Zelikman2022STaR} iteratively leverages a small set of rationale examples along with a large dataset without rationales. DAPO \citep{yu2025DAPO}, RLOO \citep{rloo}, Dr. GRPO \citep{liu2025DRGRPO}, and Open Reasoner Zero \citep{hu2025openreasonerzero} all demonstrate that minimalist, reproducible RL techniques—featuring decoupled clipping, unbiased optimization, and simple reward schemes—can significantly enhance LLM reasoning performance. 

% \noindent\textbf{Existing agent frameworks.}~~
% LLM-based agent architectures have evolved from early reasoning-action frameworks \citep{yao2022react, shinn2024reflexion, xu2023rewoo, lin2024swiftsage, wang2024mint} to structured planning approaches \citep{rafa, liu2023llm+, hao2023reasoning}. Multi-agent systems \citep{du2023improving, li2023camel, chen2023agentverse, wang2024rethinking} are designed for tasks with more complex interactions. Widely used platforms such as OpenAI Gym \citep{brockman2016openaigym} and specialized environments including Sokoban \citep{sokoban}, FrozenLake \citep{frozenlake}, and Webshop \citep{yao2022webshop} provide diverse testbeds for evaluating these agents. Moreover, general-purpose systems like HuggingGPT \citep{shen2024hugginggpt} and other frameworks \citep{wu2023autogen, hao2023toolkengpt, zhuang2023toolchain, xie2023openagents} have enabled broad applications ranging from web navigation \citep{qi2025webrltrainingllmweb}, coding copilot \citep{jimenez2024swebench,deepseekai2024deepseekcoderv2breakingbarrierclosedsource,wang2024mintevaluatingllmsmultiturn} to embodied tasks \citep{lin2024advancesembodiednavigationusing, xi2024teachingembodiedreinforcementlearning,li2025embodiedagentinterfacebenchmarking}. Social interaction capabilities have been advanced through Generative Agents and AgentSims \citep{park2023generativeagentsinteractivesimulacra,lin2023agentsimsopensourcesandboxlarge}. Challenges persist in architectural complexity and self-correction \citep{he2025enhancingllmreasoning}, especially for diverse, multi-step reasoning tasks \citep{wang2025enigmaevalbenchmarklongmultimodal,nguyen2024turningheatminpsampling, song2024trial}.


\section{Related Work}
Recent work has explored fine-tuning LLMs for multi-step reasoning through \textbf{reinforcement learning (RL)}, and developing agent frameworks to structure decision-making tasks. For reasoning, techniques range from classical PPO~\citep{schulman2017proximalpolicyoptimizationalgorithms} and actor-critic methods~\citep{haarnoja2018softactorcritic} to structured prompting via meta tokens~\citep{goyal2024thinkspeaktraininglanguage, herel2024thinkingtokenslanguagemodeling}. Policy variants such as RLOO~\citep{rloo}, GRPO~\citep{deepseekai2025deepseekr1incentivizingreasoningcapability} and DAPO~\citep{yu2025DAPO} could stabilize training and improve sample efficiency. Parallel efforts like STaR~\citep{Zelikman2022STaR} and MCTS-based reasoning~\citep{hao2023reasoning} promote step-by-step reasoning with minimal supervision. 

\textbf{On the agent side}, systems have evolved from early reactive planning~\citep{yao2022react, xu2023rewoo} to modularized decision pipelines~\citep{liu2023llm+, wu2023autogen}, multi-agent cooperation~\citep{li2023camel, wang2024rethinking}, and embodied interaction~\citep{lin2024advancesembodiednavigationusing, li2025embodiedagentinterfacebenchmarking}. Benchmarks like Sokoban~\citep{sokoban}, FrozenLake~\citep{frozenlake}, and WebShop~\citep{yao2022webshop} provide controlled testbeds for evaluating reasoning under different dynamics. Our work builds on these advances, aiming to unify RL-based reasoning with structured agent training across symbolic and language-centric tasks. We summarize more related works in Appendix~\ref{app:extendedrelatedwork}.
\vspace{-5pt}


% LIB
% Imagine while Reasoning in Space: Multimodal Visualization-of-Thought https://arxiv.org/pdf/2501.07542
% !!!!!!!!!!!!!!!!
% !!!!!!!!!!!!!!!!!!!
% !!!!!!!!!!!!!!!!!!
% Yuchen Lin Paper!!!!!!!!!!!!!!!!!!!!!!!!!!

% https://aima.cs.berkeley.edu/4th-ed/pdfs/newchap02.pdf
% https://x.com/burkov/status/1889776242573947130?s=46&t=bO8lnL2NGX4-EQlo-Qwdbw

% \paragraph{Teaching LLMs to think.}
% \paragraph{Reinforcement Learning for LLMs.}
% Reinforcement learning (RL) has significantly improved LLMs' reasoning capabilities \citep{christiano2023deepreinforcementlearninghuman,ouyang2022traininglanguagemodelsfollow}. Notable approaches include the use of Group Relative Policy Optimization (GRPO) \citep{deepseekai2025deepseekr1incentivizingreasoningcapability} for enhancing the ability of systematic problem-solving and meta tokens \citep{goyal2024thinkspeaktraininglanguage,herel2024thinkingtokenslanguagemodeling,pfau2024letsthinkdotdot} for structured thinking. DeepSeek-R1(-Zero) has demonstrated particularly strong performance by combining GRPO with meta tokens.

%\paragraph{Reinforcement Learning for LLMs.}
%Reinforcement learning (RL) has shown significant potential in enhancing LLMs' reasoning capabilities \citep{christiano2023deepreinforcementlearninghuman,ouyang2022traininglanguagemodelsfollow}.
% from basic prompting techniques \citep{reac} to more sophisticated approaches.  
% Studies revealed that increased reasoning time correlates with improved performance \citep{}, which is referred to \textit{test-time scaling law}. -> 这个和我们的work很有关系，但是这个TTS更像是training的时候很短，test的时候很长
%Subsequent work explored 
%reinforcement learning methods, including 
%Process Reward Model (PRM) \citep{zhang2025lessonsdevelopingprocessreward} and MCTS-based approaches \citep{hao2023reasoning} for systematic problem-solving. DeepSeekMath and DeepSeek-R1(-Zero) \citep{shao2024deepseekmathpushinglimitsmathematical} advanced this direction using Group Relative Policy Optimization (GRPO). Parallel developments involve meta tokens for structured thinking \citep{goyal2024thinkspeaktraininglanguage,herel2024thinkingtokenslanguagemodeling,pfau2024letsthinkdotdot} and continuous thinking tokens for reasoning compression \citep{hao2024traininglargelanguagemodels}. Among these, DeepSeek-R1(-Zero)  has demonstrated particularly strong results by combining GRPO with meta tokens.

% Teaching Large Language Models (LLMs) to think effectively has emerged as a crucial direction for reasoning tasks. Early research established that prompting techniques could effectively guide model reasoning \citep{}, leading to the discovery of a test-time compute scaling law: as models spend more time thinking through prompts, their performance improves and errors decrease \citep{}. Building on this foundation, researchers explored reinforcement learning approaches for developing systematic reasoning strategies. These include PRM-based methods \citep{}, which reward intermediate reasoning steps, and MCTS-based approaches \citep{}, which decompose problems into smaller components for systematic solution space exploration. DeepSeekMath and DeepSeek-R1(-Zero) \citep{} employ Group Relative Policy Optimization (GRPO), which innovatively estimates baselines from group scores rather than using a separate critic model. Recent advances have introduced more sophisticated approaches. DeepSeekMath \citep{} employs Group Relative Policy Optimization(GPRO), which estimates baselines from group scores rather than using a separate critic model. Another promising direction involves meta tokens, such as <think></think>, which explicitly structure the thinking process \citep{}. Parallel efforts have explored continuous thinking tokens to compress the reasoning process \citep{}. Notably, DeepSeek-R1(-Zero)  \citep{}has achieved remarkable results by combining GPRO with meta tokens, demonstrating that models can effectively develop reasoning capabilities by following this structured approach.

% \paragraph{Existing agent frameworks.}
% LLM-based agent architectures have evolved from early reasoning-action frameworks \citep{yao2022react, shinn2024reflexion, xu2023rewoo, lin2024swiftsage} to structured planning approaches \citep{liu2023llm+, hao2023reasoning}. Multi-agent systems \citep{du2023improving, li2023camel, chen2023agentverse, wang2024rethinking} are designed for tasks with more complex interactions. Challenges persist in architectural complexity and self-correction \citep{he2025enhancingllmreasoningmultipath}, especially for diverse, multi-step reasoning tasks \citep{wang2025enigmaevalbenchmarklongmultimodal,nguyen2024turningheatminpsampling, song2024trial}.
%[zihan: add trial and error paper from yuchen jin]
%\paragraph{Existing agent frameworks.}
%LLM-based agents have evolved through several stages. Initial works established reasoning-action frameworks \citep{yao2022react, shinn2024reflexion, xu2023rewoo, lin2024swiftsage}, followed by structured planning approaches \citep{liu2023llm+, hao2023reasoning, ling2024deductive}. Multi-agent systems \citep{du2023improving, li2023camel, chen2023agentverse, liang2023encouraging, chen2023reconcile, wang2024rethinking} then emerged for more complex interactions. General-purpose systems like HuggingGPT \citep{shen2024hugginggpt} and others \citep{wu2023autogen, hao2023toolkengpt, zhuang2023toolchain, xie2023openagents} enabled broad applications from web navigation \citep{qi2025webrltrainingllmweb} to reinforcement learning \citep{xi2024teachingembodiedreinforcementlearning} and embodied tasks \citep{lin2024advancesembodiednavigationusing}. Social interaction capabilities is then advanced through Generative Agents and AgentSims \citep{park2023generativeagentsinteractivesimulacra,lin2023agentsimsopensourcesandboxlarge}. However, challenges remain in architectural complexity and self-correction capabilities \citep{he2025enhancingllmreasoningmultipath}, especially for multi-step reasoning tasks \citep{wang2025enigmaevalbenchmarklongmultimodal}.

%\paragraph{Existing agent frameworks.}
%The development of LLM-based agents has seen significant progress across multiple directions. Early work established fundamental reasoning-action frameworks \citep{yao2022react, shinn2024reflexion, xu2023rewoo, lin2024swiftsage}, while subsequent research approaches enhanced structured planning \citep{liu2023llm+, hao2023reasoning, ling2024deductive}. The field then expanded into multi-agent systems \citep{du2023improving, li2023camel, chen2023agentverse, liang2023encouraging, chen2023reconcile, wang2024rethinking}, paving the way for more complex interactions. General-purpose or tool-usage agent systems like HuggingGPT \citep{shen2024hugginggpt} and others \citep{wu2023autogen, hao2023toolkengpt, zhuang2023toolchain, xie2023openagents} have enabled diverse applications ranging from web navigation \citep{qi2025webrltrainingllmweb} to reinforcement learning \citep{xi2024teachingembodiedreinforcementlearning} and embodied environments \citep{lin2024advancesembodiednavigationusing}. Social interaction capabilities have been advanced through Generative Agents and AgentSims \citep{park2023generativeagentsinteractivesimulacra,lin2023agentsimsopensourcesandboxlarge}. Despite these advances, current frameworks face key limitations in architectural complexity and self-correction capabilities \citep{he2025enhancingllmreasoningmultipath}, particularly in multi-step reasoning tasks \citep{wang2025enigmaevalbenchmarklongmultimodal,nguyen2024turningheatminpsampling}.
%[SHORTER?].
% Recent advancements in LLMs have focused heavily on agentic research. Works like ReAct \citep{}, Reflexion \citep{}, and ReWOO \citep{} integrate chain-of-thought approaches. LLM+P \citep{}, AdaPlanner \citep{}, and others \citep{} further enhance structured reasoning. Other than these single-agent settings, multi-agent paradigms have emerged through studies like CAMEL \citep{} and collaborative frameworks \citep{}, alongside the development of versatile platforms like HuggingGPT \citep{}, AutoGen \citep{}, Agent-FLAN, and related solutions \citep{}. These advancements have enabled practical applications across web navigation \citep{}, reinforcement learning \citep{}, and embodied settings including robotics \citep{} and gaming \citep{}. At the same time, human interaction has evolved through works like Generative Agents \citep{}, AgentSims \citep{}, and other social systems \citep{}. Despite their potential for handling real-world applications, current agent frameworks face key limitations. They either require complex architectural design with impractical development cycles or depend heavily on large amounts of annotated training data. Moreover, existing LLMs struggle with reliable self-correction \citep{}, leading to error accumulation in multi-step reasoning chains, which is a critical weakness for agent frameworks.